\subsection*{§. 191.\\ Beltrami 的测地表示。}

由前几节研究的 $n$ 维双曲空间在欧几里得空间上的共形表示，可以容易地导出另一种表示；Beltrami 在研究常曲率空间时正以此为出发点。在这种表示中，双曲空间测地线的像是欧几里得空间的直线（测地线）。

考虑曲率 $K=-1$ 的 $n+1$ 维双曲空间，其线元为
\[
\tag{19}
 ds^2=\frac{dx_0^2+dx_1^2+\cdots+dx_n^2}{x_0^2}.
\]
在表示欧几里得空间中考虑超球面
\[
\tag{20}
 x_0^2+x_1^2+x_2^2+\cdots+x_n^2=1,
\]
其中心位于极限超平面上；它表示（§. 188）一个测地超曲面，即一个 $n$ 维伪球空间，其曲率 $K=-1$。

这个空间的测地线也是环境 $n+1$ 维空间的测地线，因此由第422页的方程 (13) 表示：
\[
\tag{21}
\left\{
\begin{aligned}
 b x_0&=\frac1{\cosh s},\\
 x_r&=\frac{b_r}{b^2}\tgh s+c_r \qquad (r=1,2,\ldots,n),\\
 b^2&=b_1^2+b_2^2+\cdots+b_n^2,
\end{aligned}\right.
\]
不过现在为了使测地线位于超球面 (20) 上，还须在常数之间附加下列两个关系：
\[
\tag{21*}
 \sum_{r=1}^{n} b_r c_r=0,
 \qquad \frac1{b^2}+\sum_1^n c_r^2=1.
\]
现在把 $x_1,x_2,\ldots,x_n$ 看作一个表示欧几里得空间 $S_n$ 中的正交笛卡尔坐标，该空间为 $n$ 维。我们的双曲空间 (20) 全部表示在超球面
\[
 x_1^2+x_2^2+\cdots+x_n^2\leq 1,
\]
之内，而该超球面上的点表示无穷远点。又由 (21)，测地线由 $n-1$ 个关于 $x_1,x_2,x_3,\ldots,x_n$ 的线性方程
\[
 x_r=\frac{b_r}{b^2}\tgh s+c_r,
\]
表示，即由欧几里得空间 $S_n$ 的直线表示；于是测地曲面由平面表示，而一般的每个测地子流形由一个从属线性空间表示\srcfn{1}{读者可以注意到，从原来的共形表示过渡到 Beltrami 的这个测地表示，恰好通过 §. 180 所述的 Klein 几何构造实现；这里只是把该构造推广到任意维数。}。

现在需要说明：在表示欧几里得空间 $S_n$ 中取极限超球面
\[
 x_1^2+x_2^2+\cdots+x_n^2=1,
\]
内部的两点 $(x_i)$、$(x'_i)$ 时，两个对象点的测地距离 $\delta$ 按何种法则度量。为此，在 (21) 中量取弧长 $s$，其起点为 $(x'_i)$，便有
\[
 \cosh\delta=\frac{x'_0}{x_0}=\frac{\sqrt{1-\sum_i x_i'^2}}{\sqrt{1-\sum_i x_i^2}}.
\]
由 (21)、(21*) 又有
\[
 x'_r=c_r,\qquad \sum x_r x'_r=\sum x_r'{}^2,
\]
故前式可写为
\[
 \cosh\delta=\frac{1-\sum x_r x'_r}{\sqrt{1-
 \sum x_r^2}\sqrt{1-\sum x_r'^2}}.
\]
为简便起见，令
\[
 1-\sum x_r^2=\Omega_{xx},\qquad 1-\sum x_r'^2=\Omega_{x'x'},\qquad
 1-\sum x_r x'_r=\Omega_{xx'},
\]
于是
\[
\tag{22}
 \cosh\delta=\frac{\Omega_{xx'}}{\sqrt{\Omega_{xx}\Omega_{x'x'}}}.
\]
现在考察两个像点 $(x_i)$、$(x'_i)$ 的连线；它与超球面
\[
 x_1^2+x_2^2+\cdots+x_n^2=1
\]
交于两点。这两个交点与原来的两个点组成一个交比 $M$，它依赖于 $\delta$；容易得到公式
\[
\tag{22*}
 \delta=\frac12\log M.
\]
事实上，若 $(y_i)$ 是上述交点之一，则
\[
 y_i=\frac{p x_i+q x'_i}{p+q}
\]
而为确定比 $\lambda=p/q$，有二次方程
\[
 \lambda^2\Omega_{xx}+2\lambda\Omega_{xx'}+\Omega_{x'x'}=0;
\]
两根 $\lambda_1,\lambda_2$ 之比恰为 $M$，从而
\[
 \frac12(M^{1/2}+M^{-1/2})=\frac{\lambda_1+\lambda_2}{2\sqrt{\lambda_1\lambda_2}}
 =\frac{\Omega_{xx'}}{\sqrt{\Omega_{xx}\Omega_{x'x'}}},
\]
把这个公式与 (22) 比较，正好得到 (22*)。

再注意，交比 $M$ 的值由公式
\[
\tag{23}
 M=\frac{\Omega_{xx'}+\sqrt{\Omega_{xx'}^{2}-\Omega_{xx}\Omega_{x'x'}}}
 {\Omega_{xx'}-\sqrt{\Omega_{xx'}^{2}-\Omega_{xx}\Omega_{x'x'}}}.
\]
若要把双曲空间的半径取为 $R$ 而不是1，显然应把 (22*) 的右端乘以 $R$。

因此：\emph{在伪球空间位于超球面内部的测地表示中，两点的测地距离在相差一个常因子的意义下，由下列交比的对数给出：该交比由这两点在其连线上同该直线与极限超球面的两个交点组成。}
